\documentclass[a4paper,11pt]{article}

% PACKAGES
\usepackage[utf8]{inputenc}
\usepackage[T1]{fontenc}
\usepackage[english]{babel}

\usepackage[a4paper,margin=2cm]{geometry}

\usepackage{graphicx}
\usepackage{float}
\usepackage{subcaption}

\usepackage{amsmath}
\usepackage{amssymb}

\usepackage{booktabs}
\usepackage{multirow}

\usepackage{xcolor}
\usepackage{hyperref}
\usepackage{siunitx}

\hypersetup{
    colorlinks=true,
    linkcolor=blue,
    urlcolor=blue,
    citecolor=blue
}

\usepackage{enumitem}

\setlength{\parskip}{0.5em}
\setlength{\parindent}{0pt}

% DOCUMENT
\begin{document}

% TITLE PAGE
\begin{titlepage}
    \centering

    {\Large University of Trento \par}
    \vspace{0.5cm}

    {\large Department of Information Engineering and Computer Science (DISI)\par}

    \vspace{2cm}

    {\Huge \textbf{Autonomous Software Agents Project Report}\par}

    \vspace{0.5cm}

    {\Large Asa-autobots\par}

    \vfill

    \begin{flushleft}
    \textbf{Course:} Autonomous Software Agents\\
    \textbf{Authors:} Benassi Alessandro and Filippo Lollato\\
    \textbf{Academic Year:} 2025--2026
    \end{flushleft}

    \vfill

\end{titlepage}

% --------------------------------------------------
% TABLE OF CONTENTS
% --------------------------------------------------

\tableofcontents
\newpage


% --------------------------------------------------
% INTRODUCTION
% --------------------------------------------------

\section{Introduction}
This game is called Deliveroo. Agents move on a tile grid, pick up parcels across the map, 
and carry them to delivery tiles to score points. 
Parcel rewards decay over time, so the goal is to maximize total delivered score by collecting
and dropping off parcels efficiently before they lose value. 
The agents of each team cooperate to compete against the other teams.

The aim is to develop two cooperating agents:
\begin{itemize}
\item Agent 1 — BDI physical executor: does the actual moving/picking/delivering.
\item Agent 2 — LLM cognitive coordinator: has the same function of Agent 1 but extends them with the capability 
of interpreting Admin instructions received as messages in the chat, to derive policy rules, coordinating the pair 
over P2P and answer Q\&A questions to gather extra points with respect to the simple pick up and deliver process.
\end{itemize}
Other than that PDDL planning is used as an added tool to solve complex problems/situations, such as maps with crates.
\subsection{Initial state}
Both agents connect to the socket via DjsConnect() and they receive the initial state of the game.
This last consists in the static grid map, containing width, height and the tiles, which are parsed via an enum (\textit{pseudocode}):
\begin{itemize}
    \item 0: wall 
    \item 1: spawn
    \item 2: delivery
    \item 3: pavement
    \item 5: crate (5!: crate spawn)
    \item 6: parcel (6!: parcel spawn)
    \item directional one way arrow tiles (↑, ↓, ←, →)
\end{itemize}
They also recieve their own identity and state (id, name, x, y, score) via onYou and are stored in beliefs.me.
The game parameters, such as observation distance, player capacity, parcel decay interval, 
movement duration are received separately via onConfig and are kept in beliefs.config and used for calibration. 
The most used ones are cached as dedicated fields.
Both are part of the same belief base, populated at different points, 
and serve to calibrate the agent's reasoning.



% --------------------------------------------------
% BDI AGENT (shared reasoning architecture)
% --------------------------------------------------

\section{BDI Agent}

Both agents are built on the same Belief--Desire--Intention reasoning core: each
constructs a \texttt{BeliefBase} and runs an \texttt{IntentionEngine}. The physical
executor uses it directly to move, pick up and deliver; the LLM coordinator
(Section~\ref{sec:llm}) reuses the very same machinery and extends it with
natural-language reasoning. This section therefore describes the architecture once,
as the foundation shared by the whole team.

The agent maintains a model of the world (\emph{beliefs}), evaluates which goals
are worth pursuing (\emph{desires}), commits to one (\emph{intention}), and
executes it incrementally. A control loop, driven by the game loop, runs each cycle
of the schema below at roughly \SI{60}{\hertz}(tick every 16ms).

\begin{center}
\begin{verbatim}
  perceive --> revise beliefs --> deliberate --> commit to plan --> act
                    ^                (throttled)                       |
                    +-------------------- world ---------------------- +
\end{verbatim}
\end{center}

\[
\text{Beliefs} \;\rightarrow\; \text{Goal} \;\rightarrow\; \text{Plan}
\;\rightarrow\; \text{Action} \;\rightarrow\; \text{Feedback}
\;\rightarrow\; \text{Beliefs}
\]

\subsection{Sensing and Belief Revision}

Sensing happens at each movement and whenever something new enters the observation
zone. A raw perception frame from the server reports the agent's own state, the
visible map configuration, and the nearby agents, crates and parcels --- but only
within the observation radius, so it is partial and momentary. The agent never
reasons on this raw stream directly: every frame is merged into a persistent
\texttt{BeliefBase} that combines fresh observations with previously acquired
knowledge, allowing the agent to reason about entities currently out of sight.

Revision performs three non-trivial inferences:

\begin{itemize}
  \item \textbf{Spatial memory.} Parcels and crates that leave the field of view
        are retained rather than forgotten.
  \item \textbf{Negative inference.} If a remembered object should be visible, since it falls inside 
        the current radius, but it is absent from the frame, the agent
        concludes it was collected, decayed, or moved, and prunes it.
  \item \textbf{Decay simulation.} Carried and remembered parcel rewards are aged
        locally over elapsed time, so value estimates stay accurate between
        sightings.
\end{itemize}

The belief base also holds state never delivered by perception: the agent's own
identity (\texttt{me}), the carried inventory, peers, locked targets, active
cooperative contracts, and the \texttt{policyRules} issued by the LLM coordinator.
A game loop then calls the \texttt{IntentionEngine} to reason and act on the current
beliefs, while an \texttt{onMsg} handler processes Admin prompts (for the LLM agent)
and the P2P coordination messages exchanged between the two agents.

\subsection{Navigation Graph (Adjacency)}

There is no adjacency map sent by the server: the initial state (\texttt{onMap})
only provides a flat list of tiles, each with $x$, $y$ and a type. Adjacency --- the
edges of the navigation graph describing which tile-to-tile moves are legal --- is
derived on demand from these tile codes by \texttt{MapRepresentation.js}.

Two methods implement this. \texttt{isAdjacent(from, to)} checks whether a single
step is allowed: the target tile must be exactly one step away, must be walkable,
and the move must not violate the one-way arrow gate tiles (a tile cannot be entered
against its arrow direction). \texttt{getNeighbors(pos)} returns all legal neighbors
by applying \texttt{isAdjacent} to the four cardinal directions.

Because of the one-way gates, adjacency is directed and not symmetric: a move from
$A$ to $B$ may be legal while the reverse is not. For this reason no adjacency table
is precomputed; the graph is implicit and evaluated lazily. Pathfinding
(\texttt{src/mapping/Pathfinding.js}) expands it on the fly by repeatedly calling
\texttt{getNeighbors()} during the search, automatically respecting walls and gates.

\subsection{Desires and Goal Selection}

Deliberation is throttled to stay stable and avoid redundant pathfinding: goal
re-evaluation runs every \texttt{GOAL\_EVAL\_INTERVAL} (12 ticks, about 200ms), but
fires immediately when there is no active plan or when the carried inventory
changes. When it runs, \texttt{selectBestGoal(beliefs, ...)}
(\texttt{GoalSelector.js}) reads the beliefs and returns a single goal descriptor
\texttt{\{type, targetId, x, y\}}, chosen by a strict priority cascade:

\begin{enumerate}
  \item admin direct commands (\texttt{admin\_move},
        \texttt{admin\_pickup}, \texttt{admin\_deliver});
  \item active cooperative contracts with the partner
        (\texttt{rendezvous}, \texttt{handoff}, \texttt{relay});
  \item an economic, utility-based evaluation over the perceived parcels.
\end{enumerate}

The economic evaluation is the core decision and is \emph{expiration-aware}. For
each candidate parcel it estimates the value of the entire trip using A\textsuperscript{*}
travel-time estimates, discounting the raw reward by the decay accrued over that
time and adjusting it by the coordinator's policy multipliers and bonuses:
%
\begin{equation}
  \mathrm{value} = \bigl(\mathrm{reward} - r_{\text{decay}}\cdot t_{\text{trip}}\bigr)
                   \cdot m_{\text{policy}}, \qquad
  \mathrm{utility} = \frac{\mathrm{value}}{t_{\text{trip}}}.
\end{equation}
%
Because travel time is real, a distant high-reward parcel may be outscored by a
closer modest one. When already carrying cargo, the agent additionally weighs
delivering now against detouring for one more pickup, bounded by capacity, by the
decay risk to the cargo already held, and by any required-stack-size incentive, 
which can make it worth delaying delivery to assemble a larger batch. When no
option scores positively, deliberation degrades gracefully: deliver what is held,
anchor near a relay tile, patrol the parcel spawn zones, and finally wander.

\subsection{Intention Commitment and Plans}

The agent is deliberately reluctant to abandon an intention: a newly selected goal
replaces the running one only when \texttt{shouldPreemptActivePlan(...)}
(\texttt{PreemptionManager.js}) judges the switch worthwhile, which prevents
oscillation between comparable options. The exception is a path blocked by a crate,
which suspends the active plan and injects a PDDL \texttt{clear\_corridor} plan
instead (Section~\ref{sec:pddl}).

On commit, \texttt{instantiatePlanRecipe(goal)} maps the goal type to a
generator-based plan recipe (e.g. \texttt{pickup} $\to$ \texttt{CollectAndDeliver},
\texttt{deliver} $\to$ \texttt{\_deliverRecipe}, plus the admin and cooperation
recipes). These generators yield one primitive step at a time and can be suspended
and resumed via a suspension stack, so a corridor-clearing plan can interrupt a
delivery and the delivery resumes afterwards. Suspended plans that have gone stale, since 
their target disappeared or became obsolete, are discarded rather than resumed.

\subsection{Execution}

Each cycle advances the active plan by one primitive step (\texttt{move},
\texttt{pickup}, \texttt{putdown}, \texttt{wait}, \texttt{say}), handed to
\texttt{dispatchAction} (\texttt{ActionDispatcher.js}) which emits it on the socket.
The plan is told whether the previous step succeeded, so it can react to failures
rather than continue blindly.

Two behaviours sit outside the normal deliberation cycle. The first is an
\emph{opportunistic pickup}: if the agent finds itself standing on a free parcel and
is still below its carrying capacity(starting at infinity), it grabs the parcel immediately, before any
goal evaluation(deliberation). This is done to avoid leaving reachable parcels on the floor and wasting potential value.

The second is \emph{collision avoidance} between the two teammates, designed so they
never deadlock over a tile during testing. Before moving, each agent uses the partner's broadcast
position and planned next step to detect a possible conflict where both aim for the same tile. 
When this happens one agent must yield, and this is decided deterministically by comparing agent identifiers, 
so both reach the same conclusion independently without any negotiation. 
If a collision occurs anyway, the agent first waits briefly and retries(100ms), then, if the tile keeps
rejecting it, marks that tile as blocked and re-plans a route around it. After that the blocked tile is cleared so it's reusable in the future. 
When a plan finishes, the engine resumes any suspended plan or clears the goal and
re-deliberates on the next tick.


% --------------------------------------------------
% PDDL USE
% --------------------------------------------------

\section{PDDL Use}
The BDI reasoning loop described above is complemented by a second planning layer
used for one specific sub-problem: clearing pushable crates that block a corridor.
PDDL (Planning Domain Definition Language) is an exact symbolic planner. It cannot
reason about reward economics, but it is ideal for the discrete, combinatorial
question ``what sequence of moves and pushes relocates this crate out of the way?'',
where an ad-hoc heuristic would be brittle.

\subsection{Role and scope: a macro-planner, not a pathfinder}

The deliberate design choice is to invoke the planner as narrowly as possible.
Ordinary tile-by-tile navigation is handled by A\textsuperscript{*} over the grid,
which is fast and runs constantly. The PDDL solver is reserved for the
\emph{crate-clearing macro-action} alone. The reason is cost: a general-purpose
solver invoked for every navigation query would be far too slow to drive a
\SI{60}{\hertz} control loop. Restricting it to occasional obstacle resolution
keeps the symbolic machinery off the hot path while still giving the agent a
provably correct manoeuvre when it genuinely needs one.

\subsection{Local solver on port 5001}

The planner is reached through the \texttt{pddl-client} online-solver interface,
but instead of the public hosted endpoint we point it at a \emph{local} planning
service running on \texttt{http://localhost:5001} (configured via the
\texttt{PAAS\_HOST} environment variable). The motivation is purely performance:
the public online solver takes on the order of
\SIrange{1}{3}{\second} per query, with additional network latency, whereas a local
instance returns plans far faster and without round-trip overhead. Since crate
clearing temporarily suspends the agent's active plan, shaving that latency
directly reduces the time the executor stands idle in front of an obstacle.

\subsection{The domain}

The domain, \texttt{deliveroo-crates}, is a compact STRIPS model with two families
of actions:

\begin{itemize}
  \item \textbf{Movement} (\texttt{move-right/left/up/down}): the agent steps from
        a tile to an adjacent, \texttt{clear} tile in a given direction.
  \item \textbf{Pushing} (\texttt{push-right/left/up/down}): when the agent, a
        crate, and a destination tile are collinear in one direction, the agent
        advances onto the crate's tile and the crate advances to the destination---
        provided the destination is \texttt{clear} and \texttt{crate-move-capable}.
        The effects flip the \texttt{clear} predicates accordingly.
\end{itemize}

Directionality is encoded as explicit adjacency predicates
(\texttt{right/left/up/down}) rather than coordinates, so the planner reasons over
a symbolic graph. This way the one-way gate rules are respected, since adjacency is derived only
when the legality check is satisfied.

\subsection{The problem}

Each time a crate must be cleared, a fresh problem is compiled from the agent's
current beliefs. The compiler scans the walkable tiles and produces four kinds of
facts:

\begin{itemize}
  \item \textbf{Tiles and connectivity:} one object per walkable tile, together with
        the directional adjacency facts, obtained from the same move-legality check
        the agent uses for ordinary navigation.
  \item \textbf{Crate-capable tiles:} a \texttt{crate-move-capable} flag on the tiles
        that can physically hold a crate (the crate-spawn and crate-move types).
  \item \textbf{Free tiles:} a \texttt{clear} flag on every tile not occupied by a
        crate or a peer agent (both treated as obstacles).
  \item \textbf{Objects:} the agent's position, the target crate, and every other
        crate, each as a distinct object.
\end{itemize}

A design choice that has been applied is that only \emph{physical} obstacles mark a tile
as non-clear. \emph{Transient} goal blocks, such as a delivery zone that is
momentarily unreachable, are left out, because including them would label otherwise
reachable tiles as blocked and could lead the planner to wrongly conclude that no
solution exists.

The goal is a single fact: the blocking crate relocated onto a chosen clearance
tile, that is, an empty crate-capable tile out of the agent's way. The plan the
solver returns, a sequence of moves and pushes, is then translated back into grid
steps and handed to the executor as an interruptible plan.

\subsection{Focus on crates and the tiered strategy}

Crates are the only dynamic, \emph{relocatable} obstacle in the Deliveroo game, which is why
they receive dedicated planning. The agent first picks a sensible destination for
the blocking crate: a nearby clear, crate-capable tile, reachable via a push whose
``push-from'' side the agent can actually stand on (found by a short breadth-first
search). Resolution then proceeds in tiers, cheapest first:

\begin{center}
\begin{verbatim}
  A* path blocked by crate
        |
        v
  pick destination tile  --(none)-->  block goal temporarily
        |
        v
  local push solver  --(success)-->  execute push          (0 ms)
        |
     (fails)
        v
  PDDL solver @ :5001  --(success)-->  execute plan
        |
     (fails)
        v
  record cooldown + block goal
\end{verbatim}
\end{center}

A purely local, single-push computation is tried first because it is instantaneous
and covers the common case where the crate is bypassable using and adjcent tile. 
Only when that fails, typically because the push
requires manoeuvring around further crates, the agent falls back to the
PDDL solver, which can chain multiple pushes and moves. Failed solves are recorded
with a cooldown so the agent does not repeatedly hammer the planner on an
unsolvable blockage; in that case the blocked target is shelved temporarily and the
deliberation layer simply chooses something else to do and so picks a different goal.



% --------------------------------------------------
% LLM AGENT
% --------------------------------------------------

\section{LLM Agent}
\label{sec:llm}

This second agent is a cognitive coordinator. It is built on the same BDI core as the
executor, owning a \texttt{BeliefBase} and an \texttt{IntentionEngine}, so it can
itself move, pick up, and deliver, including the PDDL-based crate clearing described
above. What sets it apart is an additional natural-language reasoning layer. Its
distinctive job is to receive instructions from the Admin in free-form chat, reason
about them, and translate them into concrete commands and policies for the whole
team. It never invents game actions from this reasoning directly: every effect is
mediated either by a tool or by a peer-to-peer message, and any movement it commands,
whether for itself or the partner, is ultimately carried out by a BDI agent through
the loop of the previous sections.

\begin{center}
\begin{verbatim}
  Admin prompt --> reasoning loop --> tool call --> {local effect, P2P message}
                        ^                  |
                        +---- tool result -+        (repeat until "stop")
\end{verbatim}
\end{center}

\subsection{The reasoning loop}

A prompt triggers an iterative, single-step reasoning loop driven by a language
model under a fixed system prompt. The model returns one strictly validated JSON
object per turn, of exactly one of three kinds:

\begin{itemize}
  \item \textbf{tool}, a request to execute one of the available tools;
  \item \textbf{answer}, a piece of textual output to return;
  \item \textbf{stop}, a signal that the prompt has been fully handled.
\end{itemize}

Chain-of-thought is elicited through delimited reasoning that is stripped before
parsing, and the model runs at zero temperature for determinism. After each turn the
coordinator executes the requested tool (or records the answer), feeds the result
back as the next input, and repeats until a \texttt{stop} is produced. Outputs that
violate the schema are rejected and re-requested with an error description, up to a
retry bound. The loop advances one step at a time by design, so a prompt bundling
several tasks is decomposed and resolved in presentation order, with intermediate
answers returned in that same order.

\subsection{Prompt construction}

The system prompt is fixed and structured into a role description, explicit
definitions of \emph{reward}, \emph{task}, and \emph{rule}, the response-format
contract, the available tools, and a bank of worked examples. The definitions are
what let the model reliably tell a one-off rewarded task from a standing rule, the
distinction that drives feasibility gating. The behaviour is taught largely by
example: the prompt is few-shot, with curated cases for feasibility checks,
multi-task decomposition, rule application, and cooperation, which is what keeps a
mid-size model on the strict JSON protocol and the one-step-at-a-time discipline.
The tools section of the prompt is generated automatically from the same registry
the coordinator executes against, so the descriptions the model reads can never
drift from the tools that actually run.

\subsection{Message filtering and history}

Two filters precede the reasoning loop, both enforced in code rather than left to the
prompt. First, only messages originating from the Admin are treated as candidate
instructions; messages from any other agent are routed to the P2P handler instead.
Second, the coordinator distinguishes actionable instructions from Q\&A messages
carrying no reward: a rewardless, non-command message produces a \texttt{stop} with
no chat output at all. Conversation history is kept in full but passed to the model
only when needed and in filtered form, namely past questions and answers only, never
the intermediate reasoning or tool calls, so the prompt stays compact.

\subsection{Feasibility gating}

The central reasoning discipline is that effectful actions are gated on a feasibility
check. Any expression, whether a reward value or a logical condition such as a parity
test, must be resolved through the arithmetic tool rather than computed by the model
itself. World- or rule-changing actions (movement, variable setting, cooperation
proposals) are \textbf{reward-gated}: they are issued only once the arithmetic
evaluation has confirmed a strictly positive reward for the current prompt. A prompt
with no reward yields no action and no chat output.

Two classes are intentionally exempt. \emph{Standing policy rules}
(``from now on\ldots'', penalties, multipliers) are announcements of an effect, not
rewarded tasks, so they are always applied regardless of reward, since they describe
how the team adapts to penalties. \emph{Control and query} actions (hold, resume,
context and history lookup, contract cancellation) are likewise exempt, as they carry
no reward of their own.

\subsection{Tools}

Tools fall into two groups. \emph{Read} tools retrieve facts without side effects:
arithmetic evaluation, the local context snapshot (own and peer state, parcels, map
extremes), and conversation history. \emph{Act} tools change behaviour: directing an
agent to a coordinate, applying scoring or avoidance policy rules, setting shared
variables, holding or resuming agents, and proposing cooperation contracts.

The key design point is symmetry of effect. When an act tool targets the coordinator
itself (or ``all''), it mutates the coordinator's own beliefs directly; when it
targets the executor (or ``all''), it additionally emits a peer-to-peer message so
the executor applies the same change to its belief base. The same vocabulary governs
both agents through one uniform mechanism, detailed in Section~\ref{sec:coordination}.

% ============================================================
% COORDINATION AND MESSAGES (revised)
% ============================================================

\section{Coordination and Messages}
\label{sec:coordination}

The two agents form a team, and all of their cooperation flows through the game's
chat channel as structured messages. This section describes the transport, the
message vocabulary, and the protocols built on top of it. The same machinery serves
two roles: the LLM coordinator uses it to push commands and policies to the executor,
and both agents use it to continuously share state and resolve conflicts.

\subsection{Agent Communication}

Messages are exchanged over two primitives: \texttt{emitSay} delivers a message
privately to a single recipient, while \texttt{emitShout} broadcasts to everyone. The
team prefers the private channel, a directed \texttt{emitSay} to the partner, and
falls back to a shout only if the direct send fails. This keeps the public chat
uncluttered and avoids leaking coordination to competing teams. Incoming chat is
pre-filtered: each agent ignores its own messages and anything that does not match
the cooperative JSON schema, so ordinary Admin text and opponent chatter never reach
the protocol handlers.

Underlying everything is a steady stream of \texttt{PEER\_STATUS} heartbeats sent in
both directions a few times per second. Each heartbeat carries the sender's position,
score, carried inventory, planned next step, and known crates. This is how each agent
keeps an up-to-date model of its partner even when it is out of sensing range, and
how crate sightings are shared so both agents navigate around the same obstacles. A
lightweight \texttt{PING}/\texttt{PONG} pair serves the same purpose on demand.

\subsection{Message Structure}

Every cooperative message is a JSON object with a \texttt{type} field that selects a
handler; the remaining fields are the payload. The vocabulary groups into four
families:

\begin{itemize}
  \item \textbf{Status:} \texttt{PING}, \texttt{PONG}, \texttt{PEER\_STATUS}, for
        position, cargo and obstacle sharing.
  \item \textbf{Direct commands} (coordinator to executor): \texttt{MOVE\_TO},
        \texttt{HOLD}, \texttt{RESUME}, \texttt{APPLY\_RULES}, \texttt{SET\_VARIABLE},
        and \texttt{INSTRUCT\_SAY}. Of these, only \texttt{MOVE\_TO} expects a reply:
        the executor returns a \texttt{MOVE\_TO\_ACK} once it has reached the target
        or failed to, so the coordinator knows the move is complete.
        \texttt{INSTRUCT\_SAY} is the coordinator asking the executor to relay the
        LLM's chat answer to the Admin, a redundant delivery path that makes the
        reply more likely to reach the Admin.
  \item \textbf{Contracts:} \texttt{PROPOSE\_CONTRACT}, \texttt{ACCEPT\_CONTRACT},
        \texttt{CLOSE\_CONTRACT} for the propose/accept/close lifecycle, plus
        \texttt{SIGNAL\_READY} and \texttt{RELEASE\_CARGO} as the status transitions
        inside the hand-off and relay choreography: \texttt{SIGNAL\_READY} announces
        that an agent has dropped its cargo and stepped aside so the partner may
        collect it, and \texttt{RELEASE\_CARGO} marks the cargo as released for the
        partner to take.
  \item \textbf{Target locking:} \texttt{LOCK\_TARGET}, \texttt{RELEASE\_TARGET}
        (see the note on target deconfliction below).
\end{itemize}

On receipt, a handler applies the message directly to the agent's own belief base.
This is the mechanism by which a coordinator command takes effect remotely: a
\texttt{MOVE\_TO} becomes a high-priority \texttt{admin\_move} contract (and
implicitly releases any standing hold), an \texttt{APPLY\_RULES} updates the policy
set, \texttt{HOLD}/\texttt{RESUME} toggle the paused state, and a \texttt{SET\_VARIABLE}
writes shared memory. The executor thus needs no special command interpreter, since
coordination and perception update the same beliefs through the same revision path.

\subsection{Coordination Workflow}

\paragraph{Contract handshake.} Multi-agent manoeuvres are negotiated as
\emph{contracts} rather than hard-coded. One agent proposes, the other validates that
the target is reachable and accepts, both mark it active, and either may cancel it
later. We chose a propose/accept handshake over simply commanding the partner so that
a contract is never acted upon unless both agents agree it is feasible, which avoids
one agent waiting forever for a partner that cannot reach the meeting point. Three
cooperation patterns are built on this handshake.

\begin{itemize}
  \item \textbf{Meet-and-wait} (\texttt{RENDEZVOUS}, and \texttt{CLEARING} as a
        special case): both agents converge within a radius of a tile and pause. We
        decided to model ``clear this gate'' as a rendezvous rather than a separate
        mechanism, because mechanically it is the same act, bring an agent to a point
        and hold, and collapsing the two keeps the executor simple.
  \item \textbf{Hand-off} (\texttt{HANDOFF}): one agent carries its cargo to a shared
        tile, drops it, and steps aside so the other can collect it. This exists for
        the case where the carrier cannot reach the delivery zone but its partner can.
  \item \textbf{Relay} (\texttt{RELAY}): a persistent contract in which a designated
        courier repeatedly farms parcels and drops them, while the partner delivers
        them. The decision driving this pattern is strategic: the game awards a bonus
        when a parcel is picked up by one agent and delivered by another, and the
        relay is how the team deliberately manufactures that situation. A consequence
        we had to handle explicitly is the choice of drop tile. It is placed
        \emph{adjacent to} the best delivery zone but never \emph{on} it, because a
        courier putting a parcel down on a delivery tile would score the delivery
        itself and so cancel the very bonus the relay is meant to earn.
\end{itemize}

Blocking commands wait for confirmation before the coordinator proceeds: a directed
move resolves only once a \texttt{MOVE\_TO\_ACK} or a position match is observed, a
deliberate choice so the LLM's reasoning never advances on an unconfirmed physical
effect. Control actions such as \texttt{RESUME} clear outstanding contracts so no
agent is left waiting in a coordination loop.

\paragraph{Target deconfliction.} A lock-based deconfliction protocol is implemented
on the receiving side: an agent that committed to a parcel could announce a
\texttt{LOCK\_TARGET}, and on a clash the tie would be broken deterministically by
agent identifier (the lower identifier keeps the parcel, the other releases its claim
and re-deliberates), with \texttt{RELEASE\_TARGET} freeing a claim once the parcel is
taken or abandoned. In the final system, however, this protocol is left inactive: no
code path ever emits a \texttt{LOCK\_TARGET}, so the locks are never populated. We
compared explicit locking against relying on the continuous \texttt{PEER\_STATUS}
heartbeats and kept the latter. With only two cooperating agents, each already sees
the other's position, planned next step, path, and carried cargo several times per
second, which lets the two diverge naturally; the decay-aware utility and the
collision-yielding logic further push them apart. Explicit locking would add a
lock/release handshake, conflict tie-breaks, and stale-lock cleanup, and risks an
agent abandoning a good target on a transient clash, costs that outweighed the
marginal benefit for a two-agent team.


% --------------------------------------------------
% TESTS
% --------------------------------------------------

\section{Tests}
\subsection{Test Environment}
Both the agent have been tested along the building process continously, 
to find bugs, test different configurations, and improve them. 
Most of the test have been performed with the server and pddl running locally, this has been 
done following the instructions on the relative repositories:\dots


\subsection{Test Scenarios}
The decision was to test multiple times on the scenarios of Challenge 1 and Challenge 2.
In particular for the second part we checked that the instruciton messages sent by the Admin 
were satisfied by the agent and registered and try to fix eventual problems.

\subsection{Results and Limitations}
The most difficult part has been the prompring of the LLM agent, 
to have it doing the tasks correctly and usign the right tools.

One known limitaiton is that it is not always able to parse multiple revards for multiple requests
in one message, and therefore if the first reward is negative for example, it doesn't execute the instructions.
Another problem that has been noted is that on multiple questions, more than 
3 in one message, if there is an evaluate\_math\_expression call the model puts its output at the end.
It was difficult to spot the possible pattern that the model extracts from the prompt
in order to convince itself that the correct way to output is this one. So this 2 problems are still open 
and mostly related to the prompt.

% --------------------------------------------------
% LIMITATIONS
% --------------------------------------------------

%\section{Limitations}

%Discuss the limitations encountered during development and testing.

%Examples:

%\begin{itemize}
 %   \item Computational limitations
 %   \item LLM hallucinations
 %   \item Planner constraints
 %   \item Scalability issues
%\end{itemize}

% --------------------------------------------------
% CONCLUSION
% --------------------------------------------------

\section{Conclusion}

Summarize the main contributions and discuss future developments.

% --------------------------------------------------
% REFERENCES
% --------------------------------------------------

\begin{thebibliography}{9}

\bibitem{bdi}
Rao, A. S., Georgeff, M. P.
\textit{BDI Agents: From Theory to Practice}.

\bibitem{pddl}
McDermott et al.
\textit{PDDL -- The Planning Domain Definition Language}.

\bibitem{llm}
Recent literature on Large Language Models.

\end{thebibliography}

\end{document}